% Generated by scripts/blueprint_gen.py; do not edit.
\chapter{Proof of H2: the local gauge lemma}\label{chap:Z5}
\textbf{Target.} \texttt{theorem localGaugeLemma : LocalGaugeLemma E X} (\texttt{Cohesivemodules/Hypotheses.lean}, FROZEN), for every complex manifold \texttt{X} modelled on a finite-dimensional \texttt{E}. The statement: for every cohesive module \texttt{M} over \texttt{A = 𝒜\_X(X)} and \texttt{x ∈ X}, there are an open \texttt{U′ ∋ x}, a bounded complex \texttt{W} of finitely generated projective \texttt{O\_X(U′)}-modules, and an isomorphism in \texttt{Ho(P\_\{𝒜\_X(U′)\})} between \texttt{M|U′ = M.baseChange (dolbeaultRes E le\_top)} and \texttt{𝒜\_X(U′) ⊗\_\{O(U′)\} W = CohesiveModule.fromComplex (ρ := dolbeaultRhoU E U′) W …}.

\textbf{What we prove.} We prove the strong form, an actual isomorphism of cohesive modules with \texttt{W} free, as in Block's proof and in Bismut–Shen–Wei, Thm 5.3. Near \texttt{x}, in a chart and a frame, the cohesive operator is \texttt{𝕏 = ∂̄ + 𝒜}. Here \texttt{𝒜} is a smooth function with values in fibre operators, and flatness is the Maurer–Cartan equation \texttt{Σⱼ εⱼ ∂ⱼ𝒜 + 𝒜² = 0}. A gauge transformation, built one variable at a time following Donaldson–Kronheimer §2.2.2, makes \texttt{𝒜} free of forms and holomorphic. What remains is a complex of trivial holomorphic bundles, \texttt{W = (O\textasciicircum{}r, A₀)}.

\textbf{Five layers.}

\begin{enumerate}
\item \emph{Analytic core (C).} The one-variable gauge equation \texttt{∂̄\_k g = g·a} with parameters, solved in a finite-dimensional Banach algebra by an elementary Neumann series. It reuses the Z4 Cauchy transform.
\item \emph{Operator model (B).} Flat ∂̄-superconnections on an open subset of \texttt{E}, the gauge action, and the algebraic identities. All of this is pointwise algebra in \texttt{FibOp E r} together with derivative rules.
\item \emph{Induction (D).} Removing \texttt{dz̄₁, …, dz̄ₙ} one at a time.
\item \emph{Local freeness and frames (A, E).} The generator modules of \texttt{M|U′} are free near \texttt{x}, and the frame model connects \texttt{M|U′} to layer B.
\item \emph{Assembly (F).} \texttt{W}, the isomorphism in \texttt{HoP}, and \texttt{localGaugeLemma}.
\end{enumerate}
\textbf{Simplifications that drive the design.}

\begin{itemize}
\item The gauge equation is \emph{linear} (\texttt{∂̄\_k g = g a}), so no nonlinear fixed-point theory is needed.
\item With the operator formulation, holomorphy after each step follows by contracting the Maurer–Cartan equation with \texttt{ι\_i}. Block's p-degree separation argument is not needed (§1, B.4).
\item Smoothness comes from summable bounds on all derivatives of the Neumann series terms (\texttt{contDiff\_tsum}, in Mathlib). There is no Hölder or Sobolev space and no elliptic regularity; this closes the smoothness gap that all sources leave (§1, a5).
\item The global cutoff of Z4.e keeps every function globally defined. The slab property preserves holomorphy in the earlier variables.
\end{itemize}

\begin{theorem}[Z5.a — Kernel estimate]
\label{Z5.a}
\lean{Cohesive.integral_inv_norm_ball_zero, Cohesive.integrableOn_inv_norm_ball, Cohesive.integral_inv_norm_ball_le, Cohesive.norm_cauchySolve_le}
\leanok
\begin{itemize}
\item For ρ ≥ 0: ∫\textsubscript{B(0,ρ)} |t|\textsuperscript{−1} dA(t) = 2πρ.
\item For every c ∈ ℂ, t ↦ |t|\textsuperscript{−1} is integrable on B(c, ρ), and ∫\textsubscript{B(c,ρ)} |t|\textsuperscript{−1} dA(t) ≤ 3πρ for ρ ≥ 0.
\item Let E, Λ be complex normed spaces, v ∈ E, ℓ : E → ℂ continuous ℂ-linear with ℓ(v) = 1, and G : E → Λ with ‖G‖ ≤ C and G(w) ≠ 0 only where |ℓ(w) − c| ≤ ρ (ρ ≥ 0). Then ‖cauchySolve v G (z)‖ ≤ 3ρC for every z ∈ E.
\end{itemize}
\par\smallskip{\footnotesize Source: blueprint-ll.md §5.}
\end{theorem}
\begin{proof}
\leanok
\uses{Z4.c}
See the Lean proof and the cited items.
\paragraph{Blueprint notes.}
\texttt{∫\_\{B(c,ρ)\} \textbackslash{}|t\textbackslash{}|⁻¹ ≤ 3πρ} (polar coordinates), \texttt{norm\_cauchySolve\_le}

Files: \texttt{Analysis/KernelEstimate.lean}.
\end{proof}

\begin{theorem}[Z5.b — Derivatives of the solver]
\label{Z5.b}
\lean{Cohesive.pdbar, Cohesive.iteratedFDeriv_cauchySolve, Cohesive.pdbar_cauchySolve, Cohesive.pdbar_cauchySolve_comm}
\leanok
For complex normed spaces E, Λ, v ∈ E and f : E → Λ let ∂̄\textsubscript{v} f(z) = ½ (Df(z)(v) + i·Df(z)(i v)) (junk where f is not differentiable). Let E be finite-dimensional, Λ a Banach space, v ≠ 0 and F : E → Λ with compact support.

\begin{itemize}
\item If F is real C\textsuperscript{∞}: D\textsuperscript{m}(cauchySolve v F)(z) = cauchySolve v (D\textsuperscript{m} F)(z) for all m ∈ ℕ, z ∈ E.
\item If F is real C¹: ∂̄\textsubscript{v}(cauchySolve v F) = F, and ∂̄\textsubscript{w}(cauchySolve v F) = cauchySolve v (∂̄\textsubscript{w} F) for every w ∈ E.
\end{itemize}
\par\smallskip{\footnotesize Source: blueprint-ll.md §5.}
\end{theorem}
\begin{proof}
\leanok
\uses{Z4.c}
See the Lean proof and the cited items.
\paragraph{Blueprint notes.}
\texttt{pdbar} (+ link to \texttt{dbarPart}), \texttt{iteratedFDeriv\_cauchySolve}, \texttt{pdbar\_cauchySolve}, \texttt{pdbar\_cauchySolve\_comm} for any normed codomain; make Z4.e's \texttt{cauchySolve\_eq\_zero}, \texttt{fderiv\_cutoff\_smul*} public

Files: \texttt{Analysis/CauchySolveDeriv.lean}.
\end{proof}

\begin{theorem}[Z5.c — One-variable gauge solver]
\label{Z5.c}
\lean{Cohesive.exists_gauge_solution, Cohesive.coordCLM, Cohesive.exists_cutoff_radius, Cohesive.neumannTerm, Cohesive.gaugeSeries, Cohesive.norm_iteratedFDeriv_neumannTerm_le, Cohesive.contDiff_gaugeSeries}
\leanok
Let E be a finite-dimensional complex normed space with a fixed basis e\textsubscript{0}, …, e\textsubscript{n−1} and coordinate functionals c\textsubscript{k}, and 𝔄 a complex Banach algebra with ‖1‖ = 1.

\begin{itemize}
\item Cutoff: for a neighbourhood K of x and an index k, for all small r > 0 there is ψ as in Z4.e (with compact support in K), with 0 ≤ ψ ≤ 1 and |c\textsubscript{k}(w) − c\textsubscript{k}(x)| ≤ 2r wherever ψ(w) ≠ 0.
\item Neumann series: u\textsubscript{0} = 1, u\textsubscript{i+1} = cauchySolve v (u\textsubscript{i} · b), g = Σ\textsubscript{i} u\textsubscript{i}. If v ≠ 0, b is C\textsuperscript{∞} with compact support, ℓ(v) = 1, b vanishes outside \{|ℓ(w) − c| ≤ ρ\}, ‖b‖ ≤ B and θ = 3ρB < 1, then ‖D\textsuperscript{m} u\textsubscript{i}(z)‖ ≤ C\textsubscript{m} (i+1)\textsuperscript{m} θ\textsuperscript{i} for all m, i, z, and g is C\textsuperscript{∞} on E.
\item Main result (𝔄 finite-dimensional): let U ⊆ E be open, x ∈ U, a : E → 𝔄 real C\textsuperscript{∞} on U with values in a subalgebra S on U, and J a set of indices with k ∉ J and ∂̄\textsubscript{e\_j} a = 0 on U for j ∈ J. Then there are an open W with x ∈ W ⊆ U and g : E → 𝔄, real C\textsuperscript{∞}, with values in S and pointwise invertible, such that ∂̄\textsubscript{e\_k} g = g·a and ∂̄\textsubscript{e\_j} g = 0 (j ∈ J) on W.
\end{itemize}
\par\smallskip{\footnotesize Source: blueprint-ll.md §5.}
\end{theorem}
\begin{proof}
\leanok
\uses{Z4.e, Z5.a, Z5.b}
\textbf{C.0 Cutoff with radius control.} This is a variant of Z4.e \texttt{exists\_cutoff} exposing its construction \texttt{ψ(z) = χ₁(c\_k(z−x))·χ₂(P(z−x))}. Fix a compact neighbourhood \texttt{K ⊆ U} of \texttt{x} and \texttt{M\_K = sup\_K ‖a‖}. For every small \texttt{r > 0} there is \texttt{ψ} with:

\begin{itemize}
\item \texttt{0 ≤ ψ ≤ 1} and \texttt{tsupport ψ ⊆ K ∩ \{|c\_k(w−x)| ≤ 2r\}};
\item \texttt{ψ = 1} on an open \texttt{W ∋ x};
\item the slab \texttt{S\_r ⊇ W} is stable under \texttt{z ↦ z − t e\_k}, and \texttt{∂ⱼψ = 0} on it for \texttt{j ≠ k}.
Put \texttt{b := ψ·a} (extended by 0): \texttt{b ∈ C\_c\textasciicircum{}∞}, \texttt{‖b‖\_∞ ≤ M\_K}.

\end{itemize}
\textbf{C.3 Bounds.} Let \texttt{u₀ = 1} and \texttt{u\_\{i+1\} = T(u\_i b)}; each \texttt{u\_i} is \texttt{C\textasciicircum{}∞} with bounded derivatives (C.2, C.1).

\begin{itemize}
\item Notation: \texttt{θ = 6r M\_K}, \texttt{β\_l = sup ‖D\textasciicircum{}l b‖}, \texttt{M\_m(i) = sup ‖D\textasciicircum{}m u\_i‖}.
\item Recursion (Leibniz, \texttt{norm\_iteratedFDeriv\_mul\_le}): \texttt{M\_m(i+1) ≤ 6r Σ\_\{l≤m\} C(m,l) β\_\{m−l\} M\_l(i) ≤ θ M\_m(i) + B\_m (i+1)\textasciicircum{}\{m−1\} θ\textasciicircum{}i}.
\item Hence \texttt{M\_m(i) ≤ C\_m (i+1)\textasciicircum{}m θ\textasciicircum{}i}. Choose \texttt{r} with \texttt{θ < 1/2} first; the \texttt{β\_l} for \texttt{l ≥ 1} may then be large but are fixed. Only \texttt{β₀ ≤ M\_K} enters \texttt{θ}.
\end{itemize}
\textbf{C.4 Smoothness.} \texttt{g := Σ\_i u\_i} is \texttt{C\textasciicircum{}∞} by \texttt{contDiff\_tsum} with \texttt{v m i = C\_m (i+1)\textasciicircum{}m θ\textasciicircum{}i} (\texttt{summable\_pow\_mul\_geometric\_of\_norm\_lt\_one}).

\textbf{C.5 Equation.} \texttt{∂\_k g = Σ\_i ∂\_k u\_i = Σ\_i u\_i b = g b} (\texttt{fderiv\_tsum}, C.2, \texttt{tsum\_mul\_right}). On \texttt{W} (where \texttt{ψ = 1}), \texttt{∂\_k g = g a}.

\textbf{C.6 Subalgebra.} Let \texttt{Q : 𝔄 → 𝔄/S} be the quotient map (a CLM, finite dimension). \texttt{Q(T F) = T(Q∘F)} (\texttt{map\_cauchySolve}), so \texttt{u\_i ∈ S} by induction, and \texttt{Q(g) = Σ Q(u\_i) = 0} (\texttt{ContinuousLinearMap.map\_tsum}).

\textbf{C.7 Invertibility.} \texttt{‖g(z) − 1‖ ≤ Σ\_\{i≥1\} θ\textasciicircum{}i < 1}, so \texttt{g(z)} is a unit (\texttt{Units.oneSub}).

\textbf{C.8 Holomorphy.} For \texttt{j ∈ J}: on \texttt{S\_r}, \texttt{∂ⱼ b = (∂ⱼψ)a + ψ∂ⱼa = 0} (off \texttt{U}, \texttt{b ≡ 0} near the point). By induction \texttt{∂ⱼ u\_i = 0} on \texttt{S\_r}: \texttt{∂ⱼ u\_\{i+1\} = T(∂ⱼu\_i·b + u\_i·∂ⱼb)}, the argument vanishes on \texttt{S\_r}, and the lines \texttt{z − t e\_k} stay in \texttt{S\_r} (\texttt{cauchySolve\_eq\_zero}). Then \texttt{∂ⱼg = Σ ∂ⱼu\_i = 0} on \texttt{S\_r ⊇ W}.
\paragraph{Blueprint notes.}
cutoff with radius control (C.0), Neumann bounds (C.3), \texttt{exists\_gauge\_solution} (C.4–C.8)

Files: \texttt{Analysis/GaugeSolve.lean}.
\end{proof}

\begin{theorem}[Z5.d — Fibre operators]
\label{Z5.d}
\lean{Cohesive.dzbarTot, Cohesive.ctrTot, Cohesive.sgnTot, Cohesive.Fib, Cohesive.FibOp, Cohesive.compOp}
\leanok
Let E be a finite-dimensional complex normed space with a fixed basis e\textsubscript{0}, …, e\textsubscript{n−1}, and r ∈ ℕ. The fibre of the local model is Fib = Λ\textsuperscript{0,•}(E)\textsuperscript{r}; fibre operators are the continuous ℂ-linear endomorphisms of Fib (a finite-dimensional Banach algebra). Define: dz̄\textsubscript{j} ∈ Λ\textsuperscript{0,•}(E), the 1-forms dual to the basis; the contraction ι\textsubscript{j} of Λ\textsuperscript{0,•}(E) with e\textsubscript{j}; the grade involution σ = (−1)\textsuperscript{q} on Λ\textsuperscript{0,q}; and on Fib, componentwise, ε\textsubscript{j} = left multiplication by dz̄\textsubscript{j}, ι\textsubscript{j}, σ. Results:

\begin{itemize}
\item ε\textsubscript{iε\_j} = −ε\textsubscript{jε\_i}, ι\textsubscript{iι\_j} = −ι\textsubscript{jι\_i}, ι\textsubscript{jι\_j} = 0, ι\textsubscript{iε\_j} + ε\textsubscript{jι\_i} = δ\textsubscript{ij};
\item an operator commuting with every ε\textsubscript{j} commutes with multiplication by every t ∈ Λ\textsuperscript{0,•}(E); one anticommuting with every ε\textsubscript{j} satisfies T(t·s) = σ(t)·T(s);
\item an element of Λ\textsuperscript{0,•}(E) killed by every ι\textsubscript{j} has form degree 0.
\end{itemize}
For deg : \{0, …, r−1\} → ℤ, an element of Fib has total degree m if its i-th component has form degree m − deg(i); an operator has total degree d if it raises total degree by d.
\par\smallskip{\footnotesize Source: blueprint-ll.md §5.}
\end{theorem}
\begin{proof}
\leanok
\uses{F5}
\textbf{B.1 Coordinate formula.} For a section \texttt{s : V → Fib}, \texttt{∂̄s = Σⱼ εⱼ ∂ⱼ s}, componentwise \texttt{dbarTot}. Proof: on \texttt{Ē}, \texttt{antiPart L v = Σⱼ b\_j\textasciicircum{}*(v) antiPart L (b\_j)}, and \texttt{alternatizeUncurryFin (Σⱼ b\_j\textasciicircum{}* ⊗ cⱼ) = Σⱼ dz̄ⱼ ∧ cⱼ} (F6 wedge).
\paragraph{Blueprint notes.}
\texttt{Fib}, \texttt{FibOp}, \texttt{epsOp}, \texttt{iotaOp}, CAR relations, \texttt{IsOfTotDeg}, Λ-linearity, kernel lemma, B.1 coordinate formula

Files: \texttt{Forms/FibreOps.lean}.
\end{proof}

\begin{theorem}[Z5.e — Flat superconnections and gauges]
\label{Z5.e}
\lean{Cohesive.dbarFib, Cohesive.dbarFib_eq_sum, Cohesive.dbarFib_add_sq, Cohesive.gauge_conj, Cohesive.gaugeOp_odd, Cohesive.mcOp_gaugeOp}
\leanok
Let E be a finite-dimensional complex normed space, ∂̄\textsubscript{j} = ∂̄\textsubscript{e\_j} (Z5.b), and 𝒜, g : E → FibOp (Z5.d). For s : E → Fib, ∂̄s is taken componentwise (F9). Let mc(𝒜) = Σ\textsubscript{j} ε\textsubscript{j} ∂̄\textsubscript{j𝒜} + 𝒜² and g·𝒜 = g𝒜g\textsuperscript{−1} − (Σ\textsubscript{j} ε\textsubscript{j} ∂̄\textsubscript{j} g) g\textsuperscript{−1}. Call 𝒜 odd at z if ε\textsubscript{j} 𝒜(z) = −𝒜(z) ε\textsubscript{j} for all j.

\begin{itemize}
\item ∂̄s(z) = Σ\textsubscript{j} ε\textsubscript{j} ∂̄\textsubscript{j} s(z) if s is real-differentiable at z.
\item If 𝒜 is C¹ at z and odd at z, and s is C² at z: (∂̄ + 𝒜)((∂̄ + 𝒜)s)(z) = mc(𝒜)(z) s(z).
\item Let g be pointwise invertible and commute with every ε\textsubscript{j}. If g is C¹ at z and s differentiable at z: g(∂̄ + 𝒜)(g\textsuperscript{−1}s) = (∂̄ + g·𝒜)s at z. If 𝒜 is odd at z, so is g·𝒜. On an open U where g and 𝒜 are C\textsuperscript{∞} and 𝒜 is odd: mc(g·𝒜) = g mc(𝒜) g\textsuperscript{−1} on U.
\end{itemize}
\par\smallskip{\footnotesize Source: blueprint-ll.md §5.}
\end{theorem}
\begin{proof}
\leanok
\uses{F4, Z5.b, Z5.d}
\textbf{B.2 Curvature.} If \texttt{𝒜} is odd and Λ-linear, then \texttt{(∂̄ + 𝒜)² s = mcOp 𝒜 · s} with \texttt{mcOp 𝒜 = Σⱼ εⱼ ∂ⱼ𝒜 + 𝒜²}. This uses \texttt{∂̄² = 0} (F4), the product rule \texttt{∂ⱼ(𝒜 s) = (∂ⱼ𝒜)s + 𝒜 ∂ⱼ s}, and \texttt{𝒜εⱼ = −εⱼ𝒜}. Applied to constant sections, \texttt{𝕏² = 0 ⇔ mcOp 𝒜 = 0}.

\textbf{B.3 Gauge action.} For \texttt{g} smooth, pointwise invertible, even, Λ-linear and of degree 0, \texttt{g ∘ (∂̄ + 𝒜) ∘ g⁻¹ = ∂̄ + g·𝒜}, where \texttt{g·𝒜 := g𝒜g⁻¹ − (Σⱼ εⱼ ∂ⱼg) g⁻¹} (use \texttt{g ∂ⱼ(g⁻¹) = −(∂ⱼg)g⁻¹}).

\begin{itemize}
\item \texttt{g·𝒜} is again odd, Λ-linear and of degree 1.
\item \texttt{mcOp (g·𝒜) = g (mcOp 𝒜) g⁻¹}.
\item \texttt{(hg)·𝒜 = h·(g·𝒜)}.
\end{itemize}
\textbf{B.4 Contracted MC.} With \texttt{aᵢ := ιᵢ𝒜 + 𝒜ιᵢ} (even, degree 0): \texttt{ιᵢ mcOp 𝒜 − mcOp 𝒜 ιᵢ = ∂ᵢ𝒜 − Σⱼ εⱼ ∂ⱼaᵢ + aᵢ𝒜 − 𝒜aᵢ}.

\begin{itemize}
\item Check: \texttt{[ιᵢ, εⱼ∂ⱼ𝒜] = δᵢⱼ∂ⱼ𝒜 − εⱼ∂ⱼaᵢ} and \texttt{[ιᵢ, 𝒜²] = aᵢ𝒜 − 𝒜aᵢ}.
\item Corollary: if \texttt{mcOp 𝒜 = 0} and \texttt{aᵢ ≡ 0} on an open set, then \texttt{∂ᵢ𝒜 = 0} there.
\end{itemize}
\textbf{B.5 The step algebra.} Suppose \texttt{g} (as in B.3) commutes with \texttt{ιᵢ} for \texttt{i ∈ I ∪ \{k\}}, satisfies \texttt{∂ᵢg = 0} for \texttt{i ∈ I}, and \texttt{∂\_k g = g a\_k}. Then the coefficients of \texttt{𝒜′ = g·𝒜} are \texttt{a′ᵢ = g aᵢ g⁻¹ − (∂ᵢg)g⁻¹}; so \texttt{a′ᵢ = g aᵢ g⁻¹} for \texttt{i ∈ I}, and \texttt{a′\_k = 0}.

**B.6 Properties of \texttt{a\_k}.** If \texttt{aᵢ ≡ 0} on \texttt{U} for \texttt{i < k}, then on \texttt{U}:

\begin{itemize}
\item \texttt{a\_k} is even, Λ-linear (\texttt{[εⱼ, a\_k] = 0}) and of degree 0;
\item \texttt{[ιᵢ, a\_k] = 0} for \texttt{i ≤ k} (super Jacobi; \texttt{ι\_k² = 0});
\item \texttt{∂ᵢ a\_k = [ι\_k, ∂ᵢ𝒜]₊ = 0} for \texttt{i < k} (B.4).
\end{itemize}
\textbf{B.7 End of the induction.} If \texttt{aᵢ ≡ 0} for all \texttt{i}, then \texttt{𝒜(1 ⊗ f) ∈ Λ⁰ ⊗ F} (kernel lemma) and \texttt{𝒜(ω·s) = (−1)\textasciicircum{}\{|ω|\} ω·𝒜(s)}. So \texttt{𝒜 = σ ⊗ A₀} with \texttt{A₀ : F → F} of degree +1. By B.4, \texttt{∂ᵢA₀ = 0} for all \texttt{i}, so \texttt{A₀} is holomorphic; MC then gives \texttt{A₀² = 0}.
\paragraph{Blueprint notes.}
\texttt{mcOp}, \texttt{gaugeOp}, B.2–B.7

Files: \texttt{Forms/Superconnection.lean}.
\end{proof}

\begin{theorem}[Z5.f — Gauge induction]
\label{Z5.f}
\lean{Cohesive.gaugeSubalg, Cohesive.IsFlatSuperconn, Cohesive.IsGauge, Cohesive.exists_gauge_step, Cohesive.exists_gauge_of_aCoef, Cohesive.exists_gauge_holomorphic}
\leanok
Notation of Z5.d–Z5.e, with degrees deg : \{0, …, r−1\} → ℤ. A \emph{flat superconnection} on an open U ⊆ E is 𝒜 : E → FibOp, C\textsuperscript{∞} on U, of total degree 1, odd and with mc(𝒜) = 0 on U. A \emph{gauge} is g : E → FibOp, C\textsuperscript{∞} on E, pointwise invertible, of total degree 0, commuting with every ε\textsubscript{j}. Let a\textsubscript{i}(𝒜) = ι\textsubscript{i𝒜} + 𝒜ι\textsubscript{i} (the dz̄\textsubscript{i}-coefficient). Main result: if 𝒜 is a flat superconnection on U and x ∈ U, there are an open W with x ∈ W ⊆ U and a gauge g such that g·𝒜 is a flat superconnection on W with a\textsubscript{i}(g·𝒜) = 0 and ∂̄\textsubscript{e\_i}(g·𝒜) = 0 on W for every i. It follows by induction on k from:

\begin{itemize}
\item step: if a\textsubscript{i}(𝒜) = 0 on U for i < k, then near x there is a gauge g with a\textsubscript{i}(g·𝒜) = 0 for i ≤ k (g solves ∂̄\textsubscript{k} g = g·a\textsubscript{k} with values in the subalgebra of operators of total degree 0 commuting with every ε\textsubscript{j} and with ι\textsubscript{i} for i ≤ k; Z5.c).
\end{itemize}
\par\smallskip{\footnotesize Source: blueprint-ll.md §5.}
\end{theorem}
\begin{proof}
\leanok
\uses{Z5.c, Z5.e}
\paragraph{D. Induction (\texttt{exists\_gauge\_holomorphic})}
\begin{itemize}
\item **Claim P(k), \texttt{k = n, n−1, …, 0}.** For every open \texttt{U ∋ x} and every \texttt{𝒜} on \texttt{U} that is smooth, odd, Λ-linear, of degree 1 and MC, with \texttt{aᵢ ≡ 0} on \texttt{U} for \texttt{i < k}, there are \texttt{W ∋ x} and a gauge \texttt{g} (smooth, invertible, even, Λ-linear, degree 0) such that, on \texttt{W}, \texttt{g·𝒜 = σ ⊗ A₀} with \texttt{A₀} holomorphic and \texttt{A₀² = 0}.
\item \textbf{P(n):} take \texttt{g = 1} and B.7.
\item \textbf{P(k) from P(k+1):} apply C with:
\begin{itemize}
\item \texttt{a := a\_k};
\item \texttt{S := \{T | degree 0, [εⱼ,T] = 0 ∀j, [ιᵢ,T] = 0 ∀ i ≤ k\}};
\item \texttt{J := \{i < k\}}.
The hypotheses hold by B.6. C gives \texttt{W₁, g₁}. Then \texttt{𝒜₁ := g₁·𝒜} satisfies the hypotheses of P(k+1) on \texttt{W₁} (B.3, B.5). P(k+1) gives \texttt{W₂, g₂}; take \texttt{g := g₂g₁} (B.3 composition).

\end{itemize}
\item \texttt{exists\_gauge\_holomorphic} is P(0).
\end{itemize}
\paragraph{Blueprint notes.}
\texttt{exists\_gauge\_holomorphic} (D)

Files: \texttt{Forms/GaugeInduction.lean}.
\end{proof}

\begin{theorem}[Z5.g — Local freeness]
\label{Z5.g}
\lean{Cohesive.exists_free_baseChange, Cohesive.coef0, Cohesive.evalForm0, Cohesive.continuous_evalForm0, Cohesive.isUnit_of_evalForm0_ne_zero, Cohesive.exists_free_away, Cohesive.free_baseChange_of_isUnit}
\leanok
\begin{itemize}
\item Ring theory: a finitely generated projective module Q over a commutative ring R is free over a basic open containing a given prime 𝔪: Q[r\textsuperscript{−1}] is free over R[r\textsuperscript{−1}] for some r ∉ 𝔪. If Q[r\textsuperscript{−1}] is free and r becomes a unit in the R-algebra S, then S ⊗\textsubscript{R} Q is free over S.
\item Let E be a finite-dimensional complex normed space and X a complex manifold modelled on E. Evaluation at y ∈ U is a ring homomorphism 𝒜\textsuperscript{0,0}\textsubscript{X}(U) → ℂ; a degree-0 form is a continuous function on U, and a nowhere-vanishing one is a unit of 𝒜\textsuperscript{0,0}\textsubscript{X}(U).
\item Main result: for a cohesive module M over 𝒜\textsubscript{X}(X) with generators Q\textsuperscript{p} and x ∈ X, there is an open U′ ∋ x inside the source of chartAt(x) such that 𝒜\textsuperscript{0,0}\textsubscript{X}(U′) ⊗\_\{𝒜\textsuperscript{0,0}\textsubscript{X}(X)\} Q\textsuperscript{p} is a free 𝒜\textsuperscript{0,0}\textsubscript{X}(U′)-module for every p.
\end{itemize}
\par\smallskip{\footnotesize Source: blueprint-ll.md §5.}
\end{theorem}
\begin{proof}
\leanok
\uses{F12}
\textbf{A.1} A f.g. projective module \texttt{P} over \texttt{R = A⁰(X)} (degree-0 Dolbeault forms, i.e. smooth ℂ-valued functions) is the image of an idempotent \texttt{ε ∈ M\_N(R)}.

\textbf{A.2} \texttt{y ↦ ε(y) ∈ M\_N(ℂ)} is continuous and idempotent, so \texttt{rank ε(y) = tr ε(y)} is locally constant, equal to \texttt{r} near \texttt{x}. Choose \texttt{c₁…c\_r ∈ ℂ\textasciicircum{}N} with \texttt{ε(x)cᵢ} a basis of \texttt{im ε(x)}, and an \texttt{r × r} minor \texttt{μ(y)} of \texttt{[ε(y)c₁ … ε(y)c\_r]} with \texttt{det μ(x) ≠ 0}. On an open \texttt{U′ ∋ x}, \texttt{det μ ≠ 0} and \texttt{rank ε = r}.

\textbf{A.3} On \texttt{U′} the sections \texttt{σᵢ = ε cᵢ} form a basis of \texttt{A⁰(U′) ⊗\_R P}:

\begin{itemize}
\item they span, because \texttt{s = εs} gives \texttt{s(y) = Σ fᵢ(y) ε(y)cᵢ} with \texttt{f = μ⁻¹ ·} (the rows of \texttt{s}), which is smooth since \texttt{det μ ≠ 0} (Cramer; inverses of nonvanishing degree-0 forms are Dolbeault forms);
\item they are independent pointwise.
Shrink \texttt{U′} into \texttt{(chartAt x).source}. This gives \texttt{exists\_free\_baseChange}.

\end{itemize}
\begin{itemize}
\item \textbf{As built (blueprint §11 v1.17, recorded here by v1.18).} A.1–A.3 are replaced by localization (no idempotents, minors or Cramer): \texttt{m := ker (ev\_x : A⁰(X) → ℂ)} is prime; each \texttt{Q\textasciicircum{}p} has a free localization at some \texttt{r\_p ∉ m} (\texttt{Module.FinitePresentation.exists\_free\_localizedModule\_powers}); on \texttt{U′ := (chartAt x).source ∩ \{r\_p ≠ 0 for a ≤ p ≤ b\}} the \texttt{r\_p} are units, so \texttt{A⁰(U′) ⊗ Q\textasciicircum{}p} is free. Statement of \texttt{exists\_free\_baseChange} unchanged (\texttt{Forms/LocalFree.lean}).
\end{itemize}
\paragraph{Blueprint notes.}
A.1–A.3, \texttt{exists\_free\_baseChange}

Files: \texttt{Forms/LocalFree.lean}.
\end{proof}

\begin{theorem}[Z5.h — Frame model]
\label{Z5.h}
\lean{Cohesive.GradedFrame, Cohesive.GradedFrame.vec, Cohesive.GradedFrame.coefTot, Cohesive.GradedFrame.sum_coefTot, Cohesive.GradedFrame.coefTot_mem, Cohesive.GradedFrame.mat}
\leanok
Let A be a bounded graded-commutative dga (A1) over k and Q a graded A\textsuperscript{0}-module. A \emph{graded frame} of rank r is a degree function deg : \{0, …, r−1\} → ℤ and an A\textsuperscript{0}-basis (ẽ\textsubscript{i}) of ⨁\textsubscript{p} Q\textsuperscript{p} with ẽ\textsubscript{i} ∈ Q\textsuperscript{deg i} (coordinates of elements of Q\textsuperscript{p} vanish at frame vectors of other degrees).

\begin{itemize}
\item Every x ∈ A ⊗\textsubscript{A\textasciicircum{}0} ⨁\textsubscript{p} Q\textsuperscript{p} is uniquely Σ\textsubscript{i} x\textsubscript{i} (1 ⊗ ẽ\textsubscript{i}) with x\textsubscript{i} ∈ A; if x has total degree m, then x\textsubscript{i} ∈ A\textsuperscript{m − deg i}.
\item For a quasi-cohesive module N with a graded frame, the connection matrix ω\textsubscript{ji} ∈ A\textsuperscript{deg i + 1 − deg j} is the j-th coefficient of 𝕏(1 ⊗ ẽ\textsubscript{i}), and (𝕏x)\textsubscript{j} = d x\textsubscript{j} + Σ\textsubscript{i} σ(x\textsubscript{i}) ω\textsubscript{ji}.
\item Frame model: for N cohesive over 𝒜\textsubscript{X}(U) with a graded frame, x\textsubscript{0} ∈ U and φ = chartAt(x\textsubscript{0}), on V = φ(U ∩ φ.source) the chart expressions of the coefficients satisfy 𝕏 = ∂̄ + 𝒜 with 𝒜(z)(t)\textsubscript{j} = Σ\textsubscript{i} σ(t\textsubscript{i}) Ω\textsubscript{ji}(z) (Ω the chart expression of ω), and 𝒜 is a flat superconnection on V (Z5.f).
\end{itemize}
\par\smallskip{\footnotesize Source: blueprint-ll.md §5.}
\end{theorem}
\begin{proof}
\leanok
\uses{F10, F9, Z4.f, Z5.f}
\textbf{E.1 Frames.} On \texttt{U′ ⊆ chart source} with free generators (A), fix the frame \texttt{ẽᵢ} (degree \texttt{deg i}). Then \texttt{Tot(M|U′) = ⊕ᵢ DolbeaultAlg(U′) ẽᵢ} and \texttt{θ(ẽᵢ) = Σⱼ ωⱼᵢ ⊗ ẽⱼ} with \texttt{ωⱼᵢ} of degree \texttt{deg i + 1 − deg j}.

\textbf{E.2 Flat transport.}

\begin{itemize}
\item On \texttt{V = φ(U′)}, the chart expression of a section is \texttt{z ↦ (chartExpr sᵢ (z))ᵢ ∈ Fib}.
\item \texttt{𝒜(z)(t)\_j := Σᵢ σ(tᵢ)·chartExpr ωⱼᵢ(z)}, where \texttt{σ} is the degree sign on \texttt{TotFibre} and the product is in \texttt{TotFibre}.
\item Then \texttt{chartExpr(𝕏 s) = ∂̄(chartExpr s) + 𝒜·chartExpr s} (\texttt{chartExpr\_dbarSec\_eq}, \texttt{chartExpr\_mul}, B.1).
\item \texttt{𝒜} is smooth (F9.a), odd, Λ-linear and of degree 1; MC follows from \texttt{𝕏² = 0} via B.2, because every smooth \texttt{V → Fib} is the chart expression of a section (\texttt{chartForm}, Z4.f).
\item \emph{As built (v1.18 record):} MC is proved from \texttt{𝕏² = 0} on the frame vectors only, plus Λ-linearity of \texttt{mcOp}: \texttt{mcOp} is Λ-linear and even, so vanishing on the \texttt{r} constant generators of \texttt{Λ ⊗ ℂ\textasciicircum{}r} gives \texttt{mcOp = 0}. Same statement (\texttt{isFlatSuperconn\_frameOp}, \texttt{Forms/FrameModel.lean}); \texttt{chartForm} on all sections is not needed.
\end{itemize}
\paragraph{Blueprint notes.}
E.1–E.2 (frame, transported \texttt{𝒜}, MC from flatness)

Files: \texttt{Forms/FrameModel.lean}.
\end{proof}

\begin{theorem}[Z5.i — The local gauge lemma]
\label{Z5.i}
\lean{Cohesive.matComplex, Cohesive.exists_iso_fromComplex_of_frame, Cohesive.exists_iso_of_gauge, Cohesive.exists_local_iso}
\leanok
\begin{itemize}
\item For a commutative ring R, deg : \{0, …, r−1\} → ℤ and an r×r matrix f over R with f\textsubscript{ki} = 0 unless deg k = deg i + 1, and f² = 0: the complex W\textsuperscript{p} = R\textasciicircum{}\{\{i : deg i = p\}\}, d(w)\textsubscript{k} = Σ\textsubscript{i} f\textsubscript{ki} w\textsubscript{i}, a bounded complex of finitely generated projective R-modules.
\item Algebraic form: let A be a dga with curvature 0, ρ : R → A\textsuperscript{0} with d ∘ ρ = 0, N a cohesive module with graded frame of degrees deg and connection matrix ω, and G, H mutually inverse matrices over A with G\textsubscript{ij}, H\textsubscript{ij} ∈ A\textsuperscript{deg j − deg i} and dG\textsubscript{ij} + Σ\textsubscript{k} σ(G\textsubscript{kj}) ρ(f\textsubscript{ik}) = Σ\textsubscript{l} ω\textsubscript{lj} G\textsubscript{il}. Then N ≅ A ⊗\textsubscript{R} W in Ho(P\textsubscript{A}).
\item Geometric form: let X be a complex manifold modelled on a finite-dimensional E, x ∈ U ⊆ source of chartAt(x), N cohesive over 𝒜\textsubscript{X}(U) with a graded frame, g a gauge and B a flat superconnection on V = φ(U) with B = g·𝒜 (𝒜 the frame model), a\textsubscript{i}(B) = 0 and ∂̄\textsubscript{e\_i}B = 0 on V. Then N ≅ 𝒜\textsubscript{X}(U) ⊗\textsubscript{O\_X(U)} W in Ho(P\textsubscript{𝒜\_X(U)}), W a bounded complex of finitely generated projective O\textsubscript{X}(U)-modules.
\item Hence every cohesive module over 𝒜\textsubscript{X}(U′) with a graded frame is, near each x ∈ U′, isomorphic in Ho(P) to 𝒜\textsubscript{X}(U″) ⊗\textsubscript{O\_X(U″)} W after restriction to an open U″ ∋ x; with Z5.g this proves H2.
\end{itemize}
\par\smallskip{\footnotesize Source: blueprint-ll.md §5.}
\end{theorem}
\begin{proof}
\leanok
\uses{F12, H2, Z5.g, Z5.h}
**E.3 The complex \texttt{W}.**

\begin{itemize}
\item By D, on \texttt{W₀ ∋ φ x}: \texttt{g·𝒜 = σ ⊗ A₀}.
\item Let \texttt{U″ := chartOpen x W₀}. The entries of \texttt{A₀} are holomorphic on \texttt{W₀}; their \texttt{chartForm}s are ∂̄-closed degree-0 forms on \texttt{U″}, hence \texttt{ρ(fᵢⱼ)} for holomorphic \texttt{fᵢⱼ ∈ O(U″)} (F12).
\item Set \texttt{W\textasciicircum{}p := O(U″)\textasciicircum{}\{\#\{i : deg i = p\}\}} with differential \texttt{(fᵢⱼ)}. Then \texttt{d² = 0} (\texttt{ρ} injective, \texttt{A₀² = 0}), and \texttt{W} is bounded, f.g. and free (\texttt{IsBoundedFGProj}).
\item \emph{As built (v1.18 record):} this \texttt{ρ}-preimage construction is the one used (first option of the open question in §7; \texttt{exists\_hol\_of\_chartExpr}, \texttt{matComplex}, \texttt{Forms/LocalGauge.lean}).
\end{itemize}
\textbf{E.4 The isomorphism.} Let \texttt{G} be the \texttt{chartForm} transport of \texttt{g} (coefficients on \texttt{U″}).

\begin{itemize}
\item \texttt{Φ : fromComplex W → M|U″} via \texttt{fcExt}: \texttt{Φ(fcι wᵢ) = Σⱼ (G⁻¹)ⱼᵢ ⊗ ẽⱼ}. It is closed because \texttt{𝕏 G⁻¹ = G⁻¹(∂̄ + σ⊗A₀)} on generators, which is B.3 transported.
\item \texttt{Ψ : M|U″ → fromComplex W} on generators: \texttt{ẽⱼ ↦ Σᵢ Gᵢⱼ ⊗ fcι wᵢ} (\texttt{homEquivGen} / \texttt{famTot}-style extension).
\item \texttt{Ψ∘Φ = id} and \texttt{Φ∘Ψ = id} on generators, from \texttt{GG⁻¹ = 1} (\texttt{fc\_hom\_ext}, \texttt{hom\_ext} on generators).
\item The iso in \texttt{HoP} is built by the \texttt{DegreeZero} template. Base change from \texttt{X} to \texttt{U″} goes through \texttt{U′} (\texttt{bcCompIso}) or directly, freeness restricting along \texttt{U″ ≤ U′}.
\end{itemize}
\paragraph{F. Assembly}
\texttt{localGaugeLemma}: \texttt{intro M x}, then

\begin{enumerate}
\item A gives \texttt{U′}, frames;
\item E.1–E.2 give \texttt{𝒜} on \texttt{V};
\item D gives \texttt{W₀}, \texttt{g};
\item E.3–E.4 give \texttt{W} and the iso.
\end{enumerate}
The statement is exactly the frozen \texttt{LocalGaugeLemma E X} (\texttt{U′ := U″}).
\paragraph{Blueprint notes.}
E.3–E.4, F: \texttt{W}, \texttt{Φ}, \texttt{Ψ}, iso in \texttt{HoP}, \texttt{localGaugeLemma}

Files: \texttt{Forms/LocalGauge.lean}.
\end{proof}

\begin{theorem}[Z5.j — Main theorem]
\label{Z5.j}
\lean{Cohesive.MainStatement, main_theorem, Cohesive.blockPerf_unconditional}
\leanok
Let E be a finite-dimensional complex normed space and X a compact Hausdorff complex manifold modelled on E. Let A = 𝒜\textsubscript{X}(X) be the Dolbeault algebra of X and O\textsubscript{X} its sheaf of holomorphic functions. Then the functor J : Ho(P\textsubscript{A}) ⥤ D(O\textsubscript{X}), M ↦ α(M), into the (standard) derived category of sheaves of O\textsubscript{X}-modules is full and faithful, and its essential image is exactly D\textsubscript{perf}(X, O\textsubscript{X}). The same holds for any other \texttt{HasDerivedCategory} instance.

Deviation: D\textsubscript{perf} instead of D\textsuperscript{b\_}\{coh\}; the equivalence is stated as full, faithful and essential image D\textsubscript{perf}.
\par\smallskip{\footnotesize Source: blueprint-ll.md §5.}
\end{theorem}
\begin{proof}
\leanok
\uses{H1, H2, M, Z4.g, Z5.i}
See the Lean proof and the cited items.
\paragraph{Blueprint notes.}
Block's theorem with no hypotheses (theorem; a new frozen main statement needs your decision), \texttt{\#print axioms}

Files: \texttt{MainH2.lean}, \texttt{ChecksMain.lean}.
\end{proof}
